\documentclass[a4paper]{article}
% Español (México) con XeLaTeX: fontspec para fuentes Unicode, polyglossia
% para la división silábica y los nombres del documento en la variante mexicana.
\usepackage{fontspec}
\usepackage{polyglossia}
\setdefaultlanguage[variant=mexican]{spanish}
% Los nombres chinos van como «pinyin (original)»: xeCJK da a los caracteres
% Han su propia fuente; sin ella XeLaTeX los omite con un aviso.
% FandolSong no cubre algunos caracteres de nombres (U+73FA); WenQuanYi Zen Hei sí.
\usepackage[AutoFallBack=true]{xeCJK}
\setCJKmainfont{FandolSong-Regular.otf}
\setCJKfallbackfamilyfont{\CJKrmdefault}{WenQuanYi Zen Hei}
% polyglossia no traduce los nombres de listings.
\AtBeginDocument{\providecommand{\lstlistingname}{}\renewcommand{\lstlistingname}{Listado}%
  \providecommand{\lstlistlistingname}{}\renewcommand{\lstlistlistingname}{Índice de listados}}
\usepackage{amsmath, amssymb}
\usepackage{graphicx}
\usepackage[margin=2.35cm]{geometry}
\usepackage[most]{tcolorbox}
\usepackage{booktabs}
\usepackage{tabularx}
\usepackage{longtable}
\usepackage{array}
\usepackage{float}
\usepackage{xcolor}
\usepackage{hyperref}
\usepackage{fancyhdr}
\setCJKmainfont[AutoFakeBold=2.4]{AR PL UMing CN}
\setCJKsansfont[AutoFakeBold=2.4]{Droid Sans Fallback}
\setCJKmonofont[AutoFakeBold=2.4]{Droid Sans Fallback}
\hypersetup{colorlinks=true, linkcolor=blue!55!black, urlcolor=blue!60!black}
\graphicspath{{./}{figures/}}
\setlength{\parskip}{0.55em}
\setlength{\parindent}{2em}
\setlength{\emergencystretch}{3em}
\renewcommand{\arraystretch}{1.18}
\newtcolorbox{knowledgebox}[1]{enhanced, colback=blue!5!white, colframe=blue!70!black, colbacktitle=blue!70!black, coltitle=white, fonttitle=\bfseries, title={#1}, sharp corners, breakable}
\newtcolorbox{importantbox}[1]{enhanced, colback=yellow!10!white, colframe=yellow!75!black, colbacktitle=yellow!75!black, coltitle=black, fonttitle=\bfseries, title={#1}, sharp corners, breakable}
\newtcolorbox{warningbox}[1]{enhanced, colback=red!5!white, colframe=red!70!black, colbacktitle=red!70!black, coltitle=white, fonttitle=\bfseries, title={#1}, sharp corners, breakable}
\pagestyle{fancy}
\fancyhf{}
\fancyfoot[C]{\thepage}
\fancyfoot[R]{\footnotesize\color{gray}\href{https://github.com/hqhq1025/ai-course-notes}{AI Course Notes}}
\renewcommand{\headrulewidth}{0pt}
\renewcommand{\footrulewidth}{0pt}
\newcommand{\videourl}{https://www.youtube.com/watch?v=2o281Zy5aZE}
\newcommand{\repourl}{https://github.com/hqhq1025/ai-course-notes}

\begin{document}
\begin{titlepage}
\centering
{\Large Notas didácticas de una entrevista a fondo\par}
\vspace{1.0cm}
{\huge\bfseries Gemini Robotics, transferencia entre encarnaciones y modelos del mundo para robots\par}
\vspace{0.5cm}
{\Large Zhang Xiaojun conversa con Tan Jie (谭杰), de Google DeepMind\par}
\vspace{0.35cm}
{\large Elaborado con base en el video público de Zhang Xiaojun Podcast, una transcripción generada con GPU, la estructura de capítulos y diagramas conceptuales propios\par}
\vspace{0.3cm}
{\large 2025-11-28\par}
\vspace{0.85cm}
\includegraphics[width=0.88\textwidth,height=0.42\textheight,keepaspectratio]{cover.jpg}\par
\vfill
\begin{tcolorbox}[width=0.92\textwidth, colback=black!2!white, colframe=black!55, sharp corners]
\textbf{Título del video}: 121. Entrevista a Tan Jie (谭捷), de DeepMind: robots, transferencia entre encarnaciones, modelos del mundo, Gemini Robotics 1.5 y Google\par
\textbf{Canal del video}: Zhang Xiaojun Podcast\par
\textbf{Duración del video}: 02:06:16\par
\textbf{Enlace del video}: \href{\videourl}{\nolinkurl{\videourl}}\par
\textbf{Descripción del material}: YouTube no tiene subtítulos disponibles; el texto se elaboró con base en la transcripción de faster-whisper large-v3 con CUDA, los capítulos del video, la portada y figuras didácticas propias. Las tomas fijas de la entrevista no se reutilizan en el texto.\par
\textbf{Página del proyecto}: \href{\repourl}{\nolinkurl{\repourl}}\par
\end{tcolorbox}
\end{titlepage}

\tableofcontents
\newpage

\section{Introducción: por qué los robots están a la vez cerca y lejos}

Esta sección explica primero cómo leer este episodio. Tan Jie (谭杰) es científico investigador sénior y líder técnico del equipo de Google DeepMind Robotics. Su perspectiva sirve para separar varias preguntas que en la industria de la robótica suelen mezclarse: cómo la simulación cambia el control real, qué resolvió el aprendizaje por refuerzo, qué les aportan los grandes modelos a los robots, si los modelos fundacionales para robótica ya forman una disciplina independiente, qué resuelve el Motion Transfer de Gemini Robotics 1.5 y por qué los datos de alta calidad siguen siendo el cuello de botella más difícil.

El equilibrio más importante de este episodio es el siguiente: la robótica avanza rápido, pero todavía está lejos de implementarse a gran escala. Por un lado, Tan Jie considera que en Silicon Valley se ve a los robots, en general, como una transformación importante que está por ocurrir, y reconoce que en el último año la generalización de los modelos tuvo avances inesperados. Por otro lado, insiste en que los mejores demo que circulan en internet no suelen representar la capacidad real, y en que la implementación práctica, desde el artículo y el live demo hasta el sistema en producción, puede tardar años. Este equilibrio hace que el episodio funcione muy bien como una nota de «optimismo sereno» sobre la robótica.

\begin{importantbox}{Tesis central del episodio}
La robótica no es un problema de un solo modelo, sino un problema integral en el que intervienen a la vez el cerebro, el cerebelo, el cuerpo físico, los datos, la simulación, el tacto, la seguridad y los escenarios de aplicación. Los grandes modelos aportan comprensión del lenguaje y common sense; el aprendizaje por refuerzo resuelve parte del control de bajo nivel; los modelos del mundo y los datos sintéticos intentan cubrir la generalización; pero los datos reales o simulados de alta calidad siguen siendo el cuello de botella central.
\end{importantbox}

\begin{knowledgebox}{Criterio visual}
Este video es una entrevista con encuadre fijo, sin slides, pizarrón ni demostraciones de producto. En el texto, la portada se usa solo para identificar la fuente; todas las demás imágenes son diagramas conceptuales propios, que explican Sim2Real, VLA, Gemini Robotics 1.5, Synthetic Data, el modelo del mundo V-L-V, el tacto y la división del trabajo en robótica entre China y Estados Unidos.
\end{knowledgebox}

\subsection{Resumen de la sección}

El hilo conductor de EP121 no es que «los robots están a punto de llegar a los hogares», sino explicar por qué en la robótica ya aparecen señales de un cambio de paradigma, por qué la implementación sigue siendo lenta y qué problemas definirán el rumbo de los próximos cinco a diez años.
\section{De la graficación por computadora a la robótica: cómo el control en simulación llega al mundo real}

Esta sección comienza con la trayectoria de investigación de Tan Jie (谭杰). Su doctorado fue en graficación por computadora, en particular en physics-based character animation, es decir, lograr que figuras humanoides o animales caminen de forma natural en un entorno de simulación. Esta línea se parece mucho a la robótica: la graficación por computadora hace robótica en simulation, y la robótica hace graficación por computadora en el mundo real. La diferencia es que en la simulación se tiene acceso a toda la información de estado, mientras que el mundo real está lleno de ruido, fricción e incertidumbre.

\begin{figure}[H]
\centering
\includegraphics[width=0.96\textwidth]{graphics-to-robotics.png}
\caption{De la graficación por computadora a la robótica: hacer robótica en simulación y después transferirla al mundo real. Diagrama conceptual propio, elaborado a partir de la conversación de 00:02:00--00:13:06.}
\end{figure}

\begin{knowledgebox}{Lectura de la figura: la graficación por computadora aportó la ruta de «aprender primero en simulación»}
La figura parte de la graficación por computadora y la animación física, pasa por el aprendizaje de la marcha mediante aprendizaje por refuerzo y llega, mediante la transferencia Sim2Real, al robot real. Lo central del trabajo temprano de Tan Jie fue llevar la capacidad de control aprendida en simulación al movimiento real de robots cuadrúpedos.
\end{knowledgebox}

\subsection{Sim2Real y el aprendizaje por refuerzo: el primer cambio de paradigma}

Esta subsección aborda primero el problema de «cómo se mueve el cuerpo». Antes de la aparición de los modelos grandes, la robótica ya había vivido un giro importante: que el robot aprenda primero a moverse en simulación mediante aprendizaje por refuerzo y después transfiera lo aprendido al hardware real.

Sim2Real es la abreviatura de simulation-to-real: primero se entrena una política en un entorno de simulación y luego se transfiere al robot real. Tan Jie mencionó su trabajo temprano en Google, Sim2Real Learning Agile Locomotion for Quadruped Robots, que usó aprendizaje por refuerzo profundo para resolver la locomoción ágil de robots cuadrúpedos. El aprendizaje por refuerzo es un método de entrenamiento en el que un agente prueba y se equivoca dentro de un entorno y mejora su comportamiento a partir de la señal de reward; PPO es uno de los algoritmos más usados.

\begin{importantbox}{El primer cambio de paradigma}
Antes, el control de movimiento de los robots dependía en gran medida del control clásico y del Model Predictive Control; en los últimos cinco a diez años, el aprendizaje por refuerzo + Sim2Real transformó prácticamente el campo de la marcha y la locomotion, y volvió comunes con rapidez capacidades motrices como correr, saltar, dar volteretas o boxear.
\end{importantbox}

\subsection{El segundo cambio de paradigma, traído por los modelos grandes}

A continuación se aborda el problema de «si el robot entiende la tarea». El control de movimiento resuelve el cuerpo; los modelos grandes aportan el lenguaje, el sentido común y la planificación. Por eso Tan Jie compara los modelos grandes con el cerebro.

Antes de los modelos grandes, muchos robots carecían de common sense y tampoco entendían el lenguaje natural. Para que un robot preparara café, había que escribir un programa o descomponer la tarea en instrucciones de control muy específicas. El cambio que trajeron los modelos grandes es que el robot puede entender lenguaje natural, descomponer una tarea en pasos y saber, por sentido común, que «preparar café» requiere más o menos una taza, agua, café, calentar y servir. Tan Jie usa la analogía del cerebro y el cerebelo: el modelo grande se parece más al cerebro, encargado del pensamiento, el lenguaje y la planificación; el aprendizaje por refuerzo se parece más al cerebelo, encargado del equilibrio, el movimiento y la ejecución.

\begin{figure}[H]
\centering
\includegraphics[width=0.94\textwidth]{brain-cerebellum-robot.png}
\caption{Cerebro y cerebelo: el modelo multimodal se encarga de la cognición y el aprendizaje por refuerzo, de la ejecución. Diagrama conceptual propio, elaborado a partir de la conversación de 00:02:00--00:13:06.}
\end{figure}

\begin{knowledgebox}{Lectura de la figura: el robot necesita cerebro y también cerebelo}
El cerebro se encarga de la comprensión del lenguaje, el sentido común y la planificación de tareas; el cerebelo, de la marcha, el equilibrio, la prensión y el control de bajo nivel. Con solo cerebro, el robot «sabe qué hacer pero no lo hace bien»; con solo cerebelo, el robot «puede moverse pero no entiende el objetivo».
\end{knowledgebox}

\subsection{Resumen de la sección}

La trayectoria de investigación de Tan Jie explica los dos puntos de inflexión de la robótica en la última década: primero, el aprendizaje por refuerzo y Sim2Real transformaron el control de movimiento; segundo, los modelos grandes llevaron el lenguaje, el sentido común y la planificación a los robots. La investigación en robótica actual tiene que ocuparse a la vez del cerebro y del cerebelo.
\section{¿Es el modelo fundacional para robots una disciplina independiente? So far, not yet}

La sección anterior trató los dos cambios de paradigma; esta sección analiza una controversia: ¿debe considerarse el gran modelo fundacional para robots una disciplina independiente de los grandes modelos? La respuesta de Tan Jie (谭杰) es mesurada: so far, not yet. Hoy, la mayor parte de la inteligencia robótica sigue dependiendo de grandes modelos multimodales; lo único que se hace es complementarlos con salida de robot action, datos de acciones y fine-tuning por etapas. Si en el futuro aparecen un nuevo data format, modelos del mundo o cuellos de botella de control, podría convertirse en una disciplina más independiente, pero por ahora no ha habido un cambio cualitativo.

\begin{figure}[H]
\centering
\includegraphics[width=0.96\textwidth]{vla-action-gap.png}
\caption{VLA completa la salida de Action: para controlar un robot, un modelo multimodal tiene que aprender a producir acciones. Diagrama conceptual propio, elaborado a partir de la conversación de 00:13:06--00:23:44.}
\end{figure}

\begin{knowledgebox}{Lectura de la figura: de VLM a VLA, lo decisivo son los datos de acciones}
Vision y Language entran al modelo multimodal, que originalmente produce texto y planes; un robot necesita acciones, por lo que se incorpora Action Data para que el modelo pase de VLM a VLA. Aquí las acciones no son texto abstracto, sino articulaciones, trayectorias, efectores finales y señales de control del robot.
\end{knowledgebox}

\subsection{VLA, generalización y tasa de éxito}

Esta subsección lleva la «salida de acciones» al terreno de la utilidad práctica. La importancia de VLA no se reduce a añadir action al nombre: consiste en que el robot pase efectivamente del lenguaje y la visión a acciones ejecutables.

VLA significa Vision-Language-Action, es decir, un modelo de visión, lenguaje y acción. Busca reunir «ver el entorno», «entender la instrucción» y «producir la acción» en un mismo modelo o sistema. Los modelos actuales alcanzan tasas de éxito muy altas en tareas sencillas como pick-and-place, pero en la manipulación fina, por ejemplo subir un cierre, sujetar con precisión o controlar la orientación, la tasa de éxito puede ser de apenas 30 o 40\%. Esa cifra representa un avance en un video de investigación, pero no sirve en la vida real.

\begin{warningbox}{La tasa de éxito en investigación no equivale a la utilidad de un producto}
En el mundo real, una manipulación fina con 30\%--40\% de éxito es prácticamente inutilizable. Los usuarios necesitan sistemas estables, recuperables, explicables y que fallen de forma segura, no solo que completen la tarea una vez en un video.
\end{warningbox}

\subsection{De la idea a la demo y a la implantación}

Tan Jie recurre a una analogía con la conducción autónoma: pasar de una idea a un prototype publicado en un artículo puede tomar de seis meses a un año; del artículo a atreverse a hacer una live demo, uno o dos años; y de la live demo a la implantación real, de cinco a diez años. La robótica es más difícil que la conducción autónoma porque el espacio de acciones es mayor, hay más tareas y la interacción física es más compleja. La conducción autónoma es un escenario vertical con una salida de acciones relativamente limitada y, aun así, tardó más de diez años en acercarse a la implantación.

\begin{knowledgebox}{Tabla de etapas de la investigación a la implantación}
\begin{tabularx}{\textwidth}{p{0.22\textwidth}p{0.32\textwidth}X}
\toprule
Etapa & Producto habitual & Riesgo principal \\
\midrule
Idea & Idea de investigación, recipe de entrenamiento, hipótesis preliminar & Puede cumplirse solo en condiciones muy restringidas. \\
Prototype & Artículo, evaluación fuera de línea, video grabado & Muchos supuestos; los casos fallidos se filtran. \\
Live demo & Demostración en vivo, éxito en tareas limitadas & La estabilidad y la recuperación ante anomalías siguen siendo insuficientes. \\
Pilot & Prueba piloto con pocos clientes o escenarios & Surgen presiones de costo, mantenimiento, seguridad y operación. \\
Production & Despliegue comercial replicable & Requiere hardware confiable, un volante de datos y un sistema de servicio. \\
\bottomrule
\end{tabularx}
\end{knowledgebox}

\begin{importantbox}{La escala temporal de la robótica}
Un avance de investigación y la implantación comercial no son lo mismo. Para pasar de la demo a la producción real, un robot tiene que superar el umbral al mismo tiempo en tasa de éxito de las tareas, confiabilidad del hardware, seguridad, costo, mantenimiento, volante de datos y valor para el escenario.
\end{importantbox}

\subsection{Resumen de la sección}

Por ahora, el modelo fundacional para robots se parece más a una extensión de los grandes modelos multimodales que a una disciplina completamente independiente. Lo que de verdad decidirá si se independiza es si en el futuro hacen falta nuevos formatos de datos, modelos del mundo, paradigmas de control y mecanismos de generalización entre distintos cuerpos robóticos.
\section{El cuello de botella de los datos en robótica: el mundo real es demasiado complejo y los datos, demasiado caros}

Ya se dijo que el modelo todavía no es completamente independiente; esta sección entra en el problema más duro de la robótica: los datos. Tan Jie (谭杰) considera que el mayor problema de los robots son los datos. Los robots operan en un unstructured environment complejo, donde puede pasar cualquier cosa; las trayectorias de acción son caras, los fallos de cola larga son difíciles de cubrir y, además, las diferencias entre cuerpos de hardware dificultan compartir los datos.

\begin{figure}[H]
\centering
\includegraphics[width=0.96\textwidth]{robotics-data-bottleneck.png}
\caption{El cuello de botella de los datos en robótica: los entornos no estructurados, los fallos de cola larga y la retroalimentación de las acciones encarecen especialmente los datos. Diagrama conceptual propio, elaborado a partir de la conversación entre 00:23:44 y 00:27:52.}
\end{figure}

\begin{knowledgebox}{Lectura de la figura: los datos de robótica no son datos de video comunes}
El entorno real no es controlable, la captura de trayectorias de acción es cara, las muestras de fallo son difíciles de cubrir, entre cuerpos distintos hay diferencias de hardware y la tasa de éxito en la manipulación fina es baja. Por eso los datos de alta calidad se convierten en el cuello de botella central para que los robots pasen de la demo a la generalización.
\end{knowledgebox}

\subsection{Por qué unos cientos de horas no bastan para verificar el scaling}

Esta subsección explica por qué el problema de los datos en robótica es más difícil que el de los datos comunes de video o de texto. Entrenar un robot no consiste solo en reunir videos que «parecen relacionados», sino en reunir datos que incluyan acciones, resultados, fallos y rutas de recuperación.

Tan Jie menciona que, si en robótica solo se tienen unos cientos de horas de datos, quizá no sea posible verificar el scaling; podrían necesitarse decenas de miles de horas de datos para ver si la recipe de entrenamiento realmente funciona. Esta exigencia es muy dura tanto para los equipos de emprendimiento como para los inversionistas, porque al principio, con poco dinero y pocos datos, es muy difícil demostrar una ruta de largo plazo, y la robótica necesita justamente una convicción de largo plazo y una gran inversión.

\begin{warningbox}{Lo engañoso de las demos con pocos datos}
Con decenas o cientos de horas de datos es posible lograr una demo vistosa, pero eso no basta para demostrar generalización. Lo que la robótica necesita en realidad son datos de alta calidad que abarquen distintos escenarios, objetos, cuerpos y modos de fallo.
\end{warningbox}

\subsection{Por qué la ruta de Tesla es difícil de replicar en los robots de propósito general}

Los automóviles de Tesla son útiles, así que todos los días generan datos reales de conducción y el ciclo virtuoso de datos puede ponerse en marcha. Muchos equipos de hardware robótico pueden capturar datos, pero todavía no alcanzan el umbral de «ser útiles por sí mismos», por lo que su ciclo virtuoso de datos no arranca. Por eso no se puede afirmar sin más que «China tiene hardware y, por lo tanto, puede replicar la ruta de Tesla»: primero el hardware debe entrar en uso real para que los datos regresen de forma continua.

\subsection{Los robots son como niños pequeños: capacidades desiguales}

Esta subsección añade un juicio que se pasa por alto con facilidad: las capacidades de los robots no crecen de manera uniforme. Tan Jie recurre a la analogía con un niño: considera que la locomotion de los robots ya es muy sólida, e incluso algunas capacidades motrices superan a las de un adulto; pero la manipulation, sobre todo la manipulación con manos diestras, quizá siga al nivel de un niño de dos o tres años: entiende las instrucciones de forma aproximada, lo intenta varias veces, pero no sujeta los objetos con firmeza, y los movimientos finos le cuestan todavía más.

\begin{knowledgebox}{Capacidades desiguales: Locomotion vs Manipulation}
\begin{tabularx}{\textwidth}{p{0.24\textwidth}p{0.34\textwidth}X}
\toprule
Capacidad & Estado actual & Por qué no avanzan al mismo ritmo \\
\midrule
Locomotion & Correr, saltar, el equilibrio y la marcha ya han avanzado muy rápido & En los últimos cinco años, el aprendizaje por refuerzo y Sim2Real resolvieron en lo esencial muchos problemas de control motor. \\
Sujeción común & Puede tomar y colocar objetos sencillos y seguir instrucciones & Requiere visión, comprensión del objetivo y control grueso, pero la tolerancia a errores es relativamente alta. \\
Manipulación con manos diestras & Sigue siendo muy difícil; muchas tareas son inestables & La coordinación de varios dedos, el tacto, la fuerza de contacto y la retroalimentación del material son muy complejos. \\
Planificación cognitiva & Los modelos grandes aportan sentido común y descomposición de tareas & Pero el plan sigue dependiendo de la tasa de éxito de las acciones reales. \\
\bottomrule
\end{tabularx}
\end{knowledgebox}

\subsection{Resumen de la sección}

El problema de primeros principios de la robótica sigue siendo contar con datos de alta calidad. Sin datos suficientemente reales, diversos y verificables, no es posible juzgar si el modelo puede escalar ni pasar de la demo a un producto utilizable.
\section{Gemini Robotics 1.5: Motion Transfer y transferencia entre cuerpos robóticos}

Esta sección aborda Gemini Robotics 1.5, que da título al episodio. Tan Jie (谭杰) menciona un método clave llamado motion transfer, que consiste en transferir la capacidad de movimiento de un robot o cuerpo robótico a otro. La transferencia entre cuerpos es uno de los problemas centrales de la generalización en robótica, porque cada robot tiene distintos grados de libertad, articulaciones, manos, sensores y dinámica.

\begin{figure}[H]
\centering
\includegraphics[width=0.96\textwidth]{gemini-robotics-15-loop.png}
\caption{Ciclo cerrado de Gemini Robotics 1.5: de la visión y el lenguaje a la acción, y después generalización a nuevos cuerpos robóticos mediante Motion Transfer. Diagrama conceptual propio, elaborado a partir de la conversación de 00:27:52--00:47:32.}
\end{figure}

\begin{knowledgebox}{Lectura de la figura: Motion Transfer es el puente entre cuerpos robóticos}
Las instrucciones en lenguaje natural y el estado visual entran a Gemini, y el modelo genera planes y acciones; Motion Transfer permite que las acciones pasen de un robot a otro; al final, el robot ejecuta la acción y produce nueva retroalimentación. La dificultad está en que los cuerpos son distintos y las acciones no pueden copiarse directamente.
\end{knowledgebox}

\subsection{Por qué es difícil la transferencia entre cuerpos}

La transferencia entre cuerpos no consiste en copiar un archivo de movimiento de un robot a otro. Cada robot difiere en altura, longitud de los brazos, límites articulares, par, efector final y sensores. Al transferir el movimiento de un humanoid a un robot cuadrúpedo o a un brazo robótico, aparecen diferencias en las restricciones geométricas, dinámicas y de control. El valor de Motion Transfer consiste en abstraer la «intención de la tarea» y los «patrones de movimiento» del hardware concreto.

\begin{warningbox}{La transferencia entre cuerpos no consiste en «mezclar» los datos}
Si no se tratan las diferencias entre cuerpos, mezclar directamente las trayectorias de robots distintos puede introducir señales contradictorias. La transferencia entre cuerpos debe distinguir la intención de la tarea, la semántica del movimiento, el espacio de control y las restricciones del hardware; de lo contrario, lo que el modelo aprenda puede ser solo ruido.
\end{warningbox}

\begin{importantbox}{Generalización entre cuerpos robóticos}
La generalización entre cuerpos exige que el modelo comprenda la esencia de la tarea, no solo que memorice las trayectorias de movimiento de un robot específico. Es un paso importante para que la robótica pase del hardware especializado a capacidades generales.
\end{importantbox}

\subsection{Synthetic Data: qué hacer cuando los datos reales no bastan}

Esta sección trata una forma clave de compensar el cuello de botella de los datos. Como los datos de robots reales son caros, lentos, peligrosos y no cubren lo suficiente la cola larga, los datos sintéticos y la simulación se vuelven una opción casi ineludible.

La apuesta clave final de Tan Jie es creer en el valor de la synthetic data: con real data por sí sola no se resuelve la robótica. Synthetic data, es decir, datos sintéticos, suelen provenir de entornos de simulación, escenas generadas por programa, trayectorias automatizadas y muestreo con perturbaciones. Sus ventajas son la gran escala, la controlabilidad y la posibilidad de cubrir la cola larga; su desventaja es la sim-to-real gap, es decir, la brecha que persiste entre la simulación y la realidad.

\begin{figure}[H]
\centering
\includegraphics[width=0.96\textwidth]{synthetic-data-loop.png}
\caption{Ciclo cerrado de Synthetic Data: cuando los datos reales no bastan, la simulación genera una gran cantidad de muestras de entrenamiento controlables. Diagrama conceptual propio, elaborado a partir de 00:47:32--01:03:48 y del BAT final.}
\end{figure}

\begin{knowledgebox}{Lectura de la figura: los datos sintéticos no sustituyen a los reales, los complementan}
Las tareas reales definen el objetivo, el entorno de simulación genera escenas, las trayectorias sintéticas aportan una gran cantidad de muestras y, después del entrenamiento, hay que volver a la evaluación en el mundo real. La evaluación real expone la sim-to-real gap, que luego se retroalimenta a la simulación y al entrenamiento.
\end{knowledgebox}

\subsection{Resumen de la sección}

El valor central de Gemini Robotics 1.5 está en reunir en un mismo sistema robótico el modelo multimodal, la salida de acciones y la transferencia entre cuerpos. Motion Transfer y Synthetic Data persiguen el mismo objetivo: que las capacidades del robot se generalicen más allá de un único hardware y una única tarea.
\section{Modelos del mundo, tacto y manos diestras}

Las secciones anteriores trataron la transferencia de acciones y los datos; esta sección aborda dos capacidades más fundamentales: los modelos del mundo y el tacto. Tan Jie (谭杰) describe el modelo del mundo como Vision-Language-Vision: recibe la visión actual y el lenguaje o la acción, y genera el siguiente fotograma o el estado visual futuro. El tacto, por su parte, cobra importancia en el escenario de las manos diestras, porque el agarre fino, el deslizamiento, el material y la fuerza de contacto no pueden resolverse de forma estable solo con la visión.

\begin{figure}[H]
\centering
\includegraphics[width=0.96\textwidth]{world-model-vlv.png}
\caption{Modelo del mundo V-L-V: entrada de Vision + Language, predicción del estado visual del siguiente fotograma. Diagrama conceptual propio, elaborado a partir de la conversación en 01:03:48--01:08:29.}
\end{figure}

\begin{knowledgebox}{Lectura de la figura: el modelo del mundo «predice las consecuencias de las acciones»}
La imagen actual y el objetivo expresado en lenguaje entran al World Model; el modelo predice el siguiente fotograma o el estado visual futuro, y con base en ello la Policy elige la acción. No se limita a generar videos atractivos, sino que ayuda al robot a imaginar «si hago esto, ¿qué pasará con el mundo?».
\end{knowledgebox}

\subsection{Por qué importan los modelos del mundo}

Esta subsección lleva el modelo del mundo de la «generación de video» de vuelta al control de robots. Lo que el robot necesita es predecir las consecuencias de sus acciones, no generar una secuencia de imágenes que parezca verosímil.

El robot necesita predecir las consecuencias antes de actuar. Si al tomar una taza esta se volcará, en qué dirección hay que jalar un cierre, cómo se deformará una tela al tocarla con la mano: todo esto requiere modelar los cambios en el estado del mundo. Un modelo del mundo suficientemente capaz puede reducir el costo del ensayo y error en el mundo real y mejorar la capacidad de planificación y de recuperación.

\subsection{Por qué el tacto vuelve a ser importante}

A continuación se analizan los cambios en el nivel de los sensores. Antes se solía considerar que el tacto no era importante, en parte porque el hardware no había llegado a una etapa que exigiera un tacto fino; en cuanto se pasa a las manos diestras, el tacto deja de ser un complemento opcional y se convierte en una entrada básica.

Tan Jie menciona que, si antes se pensaba que el tacto no era importante, era por las limitaciones del hardware; con una mano diestra, el tacto se vuelve muy importante. Una pinza común puede completar muchas tareas con visión y movimientos gruesos; una mano diestra, en cambio, necesita conocer la fuerza de contacto, el deslizamiento, el material, la estabilidad del agarre y el estado de las yemas de los dedos. Sin tacto, mucha manipulation fina difícilmente se completa de forma confiable.

\begin{figure}[H]
\centering
\includegraphics[width=0.94\textwidth]{tactile-dexterous-hand.png}
\caption{Tacto y manos diestras: cuanto más diestro es el hardware, más importante es el tacto. Diagrama conceptual propio, elaborado a partir de la conversación en 01:08:29--01:17:35.}
\end{figure}

\begin{knowledgebox}{Lectura de la figura: el valor del tacto aumenta con la complejidad del cuerpo del robot}
Una pinza común se ocupa principalmente de la posición y del cierre; una mano diestra debe manejar la coordinación de varios dedos, la fuerza de contacto, el deslizamiento, el material y la estabilidad del agarre. Cuanto más se parece el hardware a la mano humana, menos es el tacto un sensor prescindible.
\end{knowledgebox}

\subsection{Resumen de la sección}

El modelo del mundo resuelve la «predicción de consecuencias» y el tacto resuelve el «contacto fino». Ambos apuntan a lo mismo: para pasar de movimientos burdos a una manipulación confiable, el robot debe comprender los cambios y los contactos del mundo físico.
\section{La ruta hacia los robots de propósito general: de la Automation a la Superhuman}

Esta sección resume cómo divide Tan Jie (谭杰) las etapas del desarrollo de los robots. Menciona cinco etapas: Automation, Teleoperation, Narrow Generalist, Home Generalist y Superhuman Capability. Esta hoja de ruta ayuda a no mezclar todas las capacidades de los robots en un solo problema.

\begin{figure}[H]
\centering
\includegraphics[width=0.96\textwidth]{robot-generalist-stages.png}
\caption{Etapas de los robots de propósito general: la ruta de largo plazo desde la automatización fija hasta las capacidades sobrehumanas. Diagrama conceptual propio, elaborado a partir de la conversación de 01:17:35--02:06:16.}
\end{figure}

\begin{knowledgebox}{Lectura de la figura: la «generalidad» es distinta en cada etapa}
La Automation consiste en reglas fijas; en la Teleoperation el hardware ya es utilizable, pero una persona lo controla a distancia; el Narrow Generalist generaliza su inteligencia dentro de un dominio acotado; el Home Generalist resuelve múltiples tareas generales en el hogar, y la Superhuman supera a los humanos en ciertas capacidades. No debe confundirse la generalización en un dominio acotado con la generalidad en el hogar.
\end{knowledgebox}

\subsection{Specialist y Generalist}

Las primeras implementaciones prácticas probablemente seguirán ocurriendo en escenarios verticales como la manufactura, la logística, los supermercados, el doblado de ropa o las servilletas. Para un specialist vertical es más fácil encontrar demanda, ROI y entornos controlables. Sin embargo, si llega a consolidarse un verdadero generalist, podrá hacer las tareas de un specialist y muchas más, y muchos specialists verán reducido su espacio. El problema es que un generalist requiere más tiempo y más datos.

\begin{warningbox}{No hay que descartar la comercialización intermedia en nombre de la forma final}
El robot humanoide de propósito general puede ser una meta de largo plazo, pero en el corto plazo los escenarios verticales todavía pueden generar valor. La ruta especializada y la ruta general no son mutuamente excluyentes; es posible que avancen en paralelo durante muchos años.
\end{warningbox}

\subsection{La seguridad no es un asunto menor}

Tan Jie afirma con claridad que hay que atender la AI safety y la robot safety. Cuando la IA o los robots puedan iterar sobre sí mismos, la humanidad enfrentará un problema de supervivencia. Google DeepMind tiene un Responsibility and Safety Council que revisa el impacto de los modelos y los robots en la sociedad, así como sus consecuencias en materia de seguridad. En el peor de los casos, si las capacidades de los robots rebasan la comprensión que se tiene de su seguridad, habría que detener la expansión de capacidades para que la investigación en seguridad se ponga al día.

\begin{importantbox}{Principios de seguridad en robótica}
Las capacidades y la seguridad deben avanzar a la par. En cada etapa del desarrollo debe hacerse la investigación de seguridad correspondiente; si el avance de las capacidades rebasa la comprensión de la seguridad, debe pausarse la expansión de capacidades para que la safety se ponga al día (catch up).
\end{importantbox}

\subsection{Por qué es fácil sobrestimar a los robots}

Esta subsección añade la advertencia de Tan Jie sobre el ánimo de la industria. Es muy fácil sobrestimar a los robots porque lo que el público ve suele ser la mejor demo: el equipo quizá la grabó diez veces y publicó la mejor toma; el video no muestra cuántas veces falló, los supuestos sobre el entorno, los reinicios manuales, los límites de la tarea ni las intervenciones de seguridad. El espectador tiende a confundir la «mejor muestra» con una «capacidad estable» y, por ello, a creer que el año próximo ya podrá comprar un robot humanoide doméstico.

\begin{warningbox}{Los límites de un video de demo como evidencia}
Un video de un robot puede demostrar que «esto ocurrió bajo ciertas condiciones», pero no que el robot tenga una capacidad desplegable. Para juzgar si está listo para usarse en la práctica hay que examinar la tasa de éxito en intentos repetidos, la recuperación ante fallas, los cambios de entorno, la diversidad de tareas, el costo, el mantenimiento y los límites de seguridad.
\end{warningbox}

La actitud de Tan Jie es de entusiasmo, pero con cabeza fría: los robots sí se están acelerando y ya empiezan a aparecer escenarios de aplicación práctica; aun así, los «robots humanoid que realmente pueden trabajar» siguen siendo hoy un desierto. Lo que de verdad hay que evitar es cometer dos errores a la vez: sobrestimar las capacidades de corto plazo y subestimar el impacto de largo plazo.

\begin{importantbox}{La percepción correcta de los plazos}
En el corto plazo, los robots todavía están lejos de la generalidad en el hogar; en el mediano plazo, es probable que primero se implementen en escenarios verticales; en el largo plazo, una vez que se consolide de verdad un generalist, el espacio de supervivencia de los specialists se reducirá.
\end{importantbox}

\subsection{Resumen de la sección}

La ruta de los robots debe verse por etapas. En el corto plazo están la generalización en dominios acotados y la implementación vertical; en el mediano plazo, el generalist doméstico, y solo en el largo plazo llegan las capacidades sobrehumanas. Los problemas de seguridad deben incorporarse al proceso de investigación desde las primeras etapas.
\section{División del trabajo en robótica entre China y Estados Unidos: hardware, convicción y volante de datos}

Esta sección vuelve a la perspectiva de China y Estados Unidos. Tan Jie (谭杰) considera que el hardware chino avanza muy rápido y que empresas como Yushu (宇树), Zhiyuan (智元) y Xinghaitu (星海图) representan las fortalezas de China en el cuerpo del robot, la cadena de suministro y la manufactura; Estados Unidos/Silicon Valley es más fuerte en la inversión a largo plazo, el cerebro del modelo, la convicción en la investigación y la capacidad de cómputo. En un escenario ideal, China y Estados Unidos deberían colaborar mejor: el desarrollo de la inteligencia en Estados Unidos y la manufactura de hardware en China se complementan.

\begin{figure}[H]
\centering
\includegraphics[width=0.94\textwidth]{china-us-robotics-split.png}
\caption{División del trabajo en robótica entre China y Estados Unidos: Estados Unidos es más fuerte en el paradigma de la inteligencia; China, en hardware y cadena de suministro. Diagrama conceptual de elaboración propia, basado en la conversación de 01:17:35--02:06:16.}
\end{figure}

\begin{knowledgebox}{Lectura de la figura: un hardware fuerte no equivale automáticamente a un volante de datos fuerte}
China es fuerte en hardware, cuerpo del robot, cadena de suministro y manufactura; Estados Unidos/Silicon Valley es fuerte en investigación a largo plazo, el cerebro del modelo, la capacidad de cómputo y la disposición a invertir en visiones de largo plazo. La robótica necesita combinar ambas cosas, pero primero el hardware debe tener un uso real para que el volante de datos empiece a girar.
\end{knowledgebox}

\subsection{La convicción de Silicon Valley y la presión de ciclos cortos en China}

Esta subsección traslada la discusión de la ruta tecnológica al entorno de recursos. La robótica requiere datos e inversión de ingeniería a largo plazo; por ello, la paciencia del capital, la capacidad de hardware y la cultura organizacional de cada región influyen directamente en que una ruta logre concretarse o no.

Tan Jie dice que Silicon Valley está dispuesto a creer en direcciones de largo plazo que parecen ambitious y que no dan resultados a corto plazo, y puede invertir dinero durante diez años; en China se tiende más a la aplicación práctica a corto plazo, la rentabilidad y el crecimiento rápido. Tanto la robótica como los modelos grandes requieren gastar mucho dinero, recolectar muchos datos y realizar experimentos de largo plazo. Si al principio solo se asignan pocos recursos, es muy difícil verificar el scaling.

\begin{warningbox}{La presión de ciclos cortos perjudica la validación de largo plazo}
Muchas rutas de la robótica no pueden validarse con el esquema de «primero dar unos pocos datos para ver el efecto y luego decidir si se invierte». Si el volumen inicial de datos es demasiado pequeño, el modelo ni siquiera puede mostrar las leyes de scaling, y el equipo podría abandonar por error una dirección de largo plazo que era correcta.
\end{warningbox}

\subsection{La ruta de Google, la ruta de Waymo y la ruta de Tesla}

Zhang Xiaojun (张小珺) propone una analogía: la ruta de robótica de Google podría parecerse a la de Waymo, con énfasis en el cerebro y el sistema; las empresas chinas de hardware podrían parecerse a Tesla, con la expectativa de que el despliegue de hardware genere datos. Tan Jie advierte que los autos de Tesla son útiles por sí mismos y por eso su volante de datos puede girar; actualmente, mucho hardware robótico aún no alcanza el umbral de utilidad, por lo que no puede compararse sin más con Tesla.

\begin{knowledgebox}{Condiciones para que funcione el volante de datos}
\begin{tabularx}{\textwidth}{p{0.24\textwidth}p{0.34\textwidth}X}
\toprule
Condición & Significado & Dificultad en robótica \\
\midrule
El producto es útil & Los usuarios están dispuestos a usarlo de verdad, y no solo a ver una demo & Muchos robots aún no alcanzan el umbral de utilidad cotidiana. \\
Escenarios de alta frecuencia & La frecuencia de uso basta para generar una gran cantidad de datos & Los entornos domésticos y abiertos tienen demasiados casos de cola larga. \\
Retroalimentación registrable & Los éxitos, los fracasos y las intervenciones humanas pueden registrarse de forma estructurada & La causa de una falla de manipulación puede estar en la visión, el control de fuerza, la planificación o el hardware. \\
Datos de retorno entrenables & Los datos pueden entrar en el ciclo cerrado de entrenamiento y evaluación & Las diferencias entre cuerpos de robot y la privacidad/seguridad aumentan el costo. \\
\bottomrule
\end{tabularx}
\end{knowledgebox}

\subsection{Resumen de la sección}

La división del trabajo en robótica entre China y Estados Unidos no consiste en que uno reemplace al otro, sino en la complementariedad entre hardware e inteligencia. El verdadero reto es llevar el hardware a escenarios útiles, formar el volante de datos y, con inversión de largo plazo, conectar los modelos, los datos sintéticos y la evaluación en condiciones reales.
\section{Términos clave: índice de palabras clave de este episodio}

\begin{longtable}{p{0.24\textwidth}p{0.34\textwidth}p{0.33\textwidth}}
\toprule
Término & Explicación breve & Papel en este episodio \\
\midrule
Robotics & Robótica: estudia la percepción, el control, la planificación y la ejecución física & Tema de todo el episodio. \\
Embodied AI & Inteligencia corporeizada: un agente percibe su entorno y actúa en él & Es un término común en el contexto chino, pero Tan Jie (谭杰) prefiere hablar de robots. \\
Sim2Real & Transferencia del entrenamiento en simulación al mundo real & Núcleo del primer cambio de paradigma de Tan Jie. \\
Reinforcement Learning & Aprendizaje por refuerzo: optimiza una política por ensayo y error mediante reward & Método importante para resolver la locomotion y el control. \\
PPO & Proximal Policy Optimization, algoritmo de aprendizaje por refuerzo de uso común & Uno de los antecedentes técnicos de los primeros trabajos de Tan Jie. \\
MPC & Model Predictive Control, control predictivo basado en modelos & Representa la línea tradicional de control de robots. \\
VLM & Vision-Language Model, modelo de visión y lenguaje & Puede ver imágenes y entender lenguaje, pero por lo general no produce acciones directamente. \\
VLA & Vision-Language-Action, modelo de visión, lenguaje y acción & Forma importante de los modelos fundacionales para robots. \\
Motion Transfer & Transferencia de movimiento: lleva las capacidades de movimiento a cuerpos robóticos distintos & Uno de los mecanismos clave de Gemini Robotics 1.5. \\
World Model & Modelo que predice cómo cambia el estado del mundo & Ayuda al robot a planificar las consecuencias de sus acciones. \\
V-L-V & Vision-Language-Vision & Una de las formas en que Tan Jie describe los modelos del mundo. \\
Synthetic Data & Datos sintéticos, generalmente obtenidos de simulación o generación por programa & Compensan la escasez de datos de robots reales. \\
Tactile Sensing & Percepción táctil: mide la fuerza de contacto, el deslizamiento y el material & Entrada importante para las manos diestras y la manipulación fina. \\
Dexterous Hand & Mano diestra: cuerpo robótico de varios dedos para operaciones complejas & Zona donde se concentran las dificultades del tacto y la manipulation. \\
Generalist & Robot de propósito general, capaz de generalizar entre tareas & Objetivo de largo plazo; terminará por desplazar al specialist. \\
Specialist & Robot especializado que resuelve tareas verticales específicas & Es más fácil de llevar a la práctica en el corto plazo. \\
\bottomrule
\end{longtable}

\subsection{Resumen de la sección}

Los términos de este episodio muestran que la robótica no es el nombre de un solo modelo, sino un conjunto de problemas de sistemas que abarcan la simulación, la acción, los modelos del mundo, el tacto, el hardware y la seguridad. Entender las relaciones entre los términos importa más que perseguir una demo aislada.
\section{Síntesis y ampliación}

\subsection{Conclusiones principales}

Esta sección condensa en una lista de verificación los juicios sobre robótica de toda la nota. El lector puede usarla para evaluar una demo de robots o el discurso de una empresa: ¿resuelve el cerebro, el cerebelo, los datos, el cuerpo del robot, el tacto y la seguridad, o solo muestra una de estas piezas?

\begin{enumerate}
\item El primer cambio de paradigma de la robótica en la última década fue que el aprendizaje por refuerzo y Sim2Real transformaron la locomotion.
\item Los modelos grandes trajeron el segundo cambio de paradigma: dieron a los robots comprensión del lenguaje, sentido común y capacidad de planificación.
\item Los modelos base para robots siguen siendo, en lo esencial, modelos grandes multimodales a los que se agregan datos de acciones y fine-tuning; todavía no son del todo independientes.
\item Los datos de alta calidad son hoy el principal cuello de botella de la robótica; unos cientos de horas de datos no bastan para verificar el scaling.
\item Motion Transfer, de Gemini Robotics 1.5, intenta resolver la generalización entre distintos cuerpos de robot.
\item Synthetic Data es una dirección clave en la que, según Tan Jie (谭杰), hay que creer: los datos reales por sí solos no alcanzan.
\item El valor de los modelos del mundo está en predecir las consecuencias de las acciones, no solo en generar video.
\item Las manos diestras vuelven a dar importancia al tacto, porque la manipulation fina necesita información de contacto.
\item Los robots de propósito general pasarán por una etapa larga que va de la automatización fija a capacidades sobrehumanas.
\item La ruta ideal para la robótica de China y Estados Unidos es la complementariedad entre el paradigma de inteligencia y la cadena de suministro de hardware.
\end{enumerate}

\subsection{Preguntas abiertas}

Esta sección conserva preguntas abiertas porque muchas de las rutas técnicas discutidas en este episodio siguen cambiando con rapidez. Lo que más necesita la industria robótica no es una respuesta aislada, sino seguir de cerca qué variables mejoran de verdad: la escala de los datos, la transferencia entre cuerpos de robot, el hardware táctil, los modelos del mundo y la evaluación de seguridad.

\begin{itemize}
\item ¿Cuándo serán los modelos base para robots realmente independientes de los modelos grandes multimodales?
\item ¿Puede Motion Transfer funcionar de manera estable en un número suficiente de cuerpos de robot?
\item ¿Cómo puede Synthetic Data reducir la sim-to-real gap?
\item En los modelos del mundo, ¿importa más la predicción visual o la predicción de acciones y estados?
\item ¿Cuándo podrá el hardware táctil incorporarse a escala al ciclo de retroalimentación de datos de las manos diestras?
\item ¿Cuánto tiempo coexistirán los robots humanoides de propósito general y los specialist verticales?
\item ¿Cómo debe la safety de los robots adelantarse al rápido crecimiento de sus capacidades?
\end{itemize}

\subsection{Lecturas adicionales}

\begin{itemize}
\item Sim2Real Learning Agile Locomotion for Quadruped Robots: para entender cómo se transfiere el aprendizaje por refuerzo a robots cuadrúpedos reales.
\item Serie Robot Transformer: trabajos como RT-1, RT-2 y RT-X.
\item Gemini Robotics / Gemini Robotics 1.5: la línea de modelos para robots de Google DeepMind.
\item EP132, entrevista con Gao Jiyang (高继扬): Xinghaitu (星海图), Waymo/Momenta y la puesta en producción de la inteligencia corporizada.
\item EP134, panorama de datos: datos para robots, Recipe y fijación de precios de los datos.
\end{itemize}

\begin{importantbox}{Juicio final}
Lo más difícil en robótica no es «lograr que el modelo entienda una frase», sino lograr que un sistema físico actúe de forma estable en el mundo real. Los modelos grandes dieron a los robots un cerebro, el aprendizaje por refuerzo aportó parte del cerebelo, y los datos sintéticos y los modelos del mundo intentan suplir la experiencia; pero la implementación real todavía requiere datos de alta calidad, hardware confiable, tacto, seguridad y paciencia a largo plazo.
\end{importantbox}

\end{document}
